\bibitem{AbSt} M. Abramowitz \& I. A. Stegun - Handbook of mathematical functions with formulas, graphs, and mathematical tables , National Bureau of Standards Applied Math. Series , vol. 55 , U.S. Government Printing Office, Washington, D.C., 1964
\bibitem{AgBaBo} A. Agrachev, D. Barilari \& U. Boscain - A comprehensive introduction to sub-Riemannian geometry , Cambridge Studies in Advanced Math. , vol. 181 , Cambridge University Press, Cambridge, 2020
\bibitem{An} C. Anné - “Bornes sur la multiplicité” (1992), Prépublications EPFL,
\bibitem{ArRi} V. Arnaiz \& G. Rivière - “Quantum limits of perturbed sub-Riemannian contact Laplacians in dimension 3” , J. Éc. polytech. Math. 11 (2024), p. 909-956
\bibitem{Ba} H. Bahouri - “Sur la propriété de prolongement unique pour les opérateurs de Hörmander” , in Journées équations aux dérivées partielles , École polytechnique, Palaiseau, 1983, article ID 15, p. 1-7
\bibitem{Ba2} H. Bahouri - “Non prolongement unique des solutions d’opérateurs ‘somme de carrés’” , Ann. Inst. Fourier (Grenoble) 36 (1986) no. 4, p. 137-155
\bibitem{BaCoLeSC} D. Bakry, T. Coulhon, M. Ledoux \& L. Saloff-Coste - “Sobolev inequalities in disguise” , Indiana Univ. Math. J. 44 (1995) no. 4, p. 1033-1074
\bibitem{BaGo} M. S. Baouendi \& C. Goulaouic - “Nonanalytic-hypoellipticity for some degenerate elliptic operators” , Bull. Amer. Math. Soc. 78 (1972), p. 483-486
\bibitem{BeRi} - Proc. conference ‘Journées non holonomes: géométrie sous-riemannienne, théorie du contrôle, robotique’ (Paris, June 30–July 1, 1992) (A. Bellaïche \& J.-J. Risler, eds.) , Progress in Math. , Birkhäuser, Basel, 1996
\bibitem{BerMe} P. Bérard \& D. Meyer - “Inégalités isopérimétriques et applications” , Ann. Sci. École Norm. Sup. (4) 15 (1982) no. 3, p. 513-541
\bibitem{Bo} J.-M. Bony - “Principe du maximum, inégalite de Harnack et unicité du problème de Cauchy pour les opérateurs elliptiques dégénérés” , Ann. Inst. Fourier (Grenoble) 19 (1969), p. 277-304
\bibitem{BrFoMo} T. P. Branson, L. Fontana \& C. Morpurgo - “Moser-Trudinger and Beckner-Onofri’s inequalities on the CR sphere” , Ann. of Math. (2) 177 (2013) no. 1, p. 1-52
\bibitem{CaDaPaTy} L. Capogna, D. Danielli, S. D. Pauls \& J. T. Tyson - An introduction to the Heisenberg group and the sub-Riemannian isoperimetric problem , Progress in Math. , vol. 259 , Birkhäuser Verlag, Basel, 2007
\bibitem{Car} G. Carron - “Inégalités isopérimétriques de Faber-Krahn et conséquences” , in Actes de la table ronde de géométrie différentielle (Luminy, 1992) , Sémin. Congr. , vol. 1 , Société Mathématique de France, Paris, 1996, p. 205-232
\bibitem{ChChHwYa} J.-H. Cheng, H.-L. Chiu, J.-F. Hwang \& P. Yang - “Umbilicity and characterization of Pansu spheres in the Heisenberg group” , J. reine angew. Math. 738 (2018), p. 203-235
\bibitem{ChHwMaYa} J.-H. Cheng, J.-F. Hwang, A. Malchiodi \& P. Yang - “A Codazzi-like equation and the singular set for C 1 smooth surfaces in the Heisenberg group” , J. reine angew. Math. 671 (2012), p. 131-198
\bibitem{ChoPo1} W. Choi \& R. Ponge - “Privileged coordinates and nilpotent approximation for Carnot manifolds, I. General results” , J. Dynam. Control Systems 25 (2019) no. 1, p. 109-157
\bibitem{ChoPo2} W. Choi \& R. Ponge - “Privileged coordinates and nilpotent approximation for Carnot manifolds, II. Carnot coordinates” , J. Dynam. Control Systems 25 (2019) no. 4, p. 631-670
\bibitem{CoLu} W. S. Cohn \& G. Lu - “Best constants for Moser-Trudinger inequalities on the Heisenberg group” , Indiana Univ. Math. J. 50 (2001) no. 4, p. 1567-1591
\bibitem{CHT0} Y. Colin de Verdière, L. Hillairet \& E. Trélat - “Spectral asymptotics for sub-Riemannian Laplacians, I: Quantum ergodicity and quantum limits in the 3-dimensional contact case” , Duke Math. J. 167 (2018) no. 1, p. 109-174
\bibitem{CHT} Y. Colin de Verdière, L. Hillairet \& E. Trélat - “Small-time asymptotics of hypoelliptic heat kernels near the diagonal, nilpotentization and related results” , Ann. H. Lebesgue 4 (2021), p. 897-971
\bibitem{CHT1} Y. Colin de Verdière, L. Hillairet \& E. Trélat - “Spectral asymptotics for sub-Riemannian Laplacians” , 2022
\bibitem{CH} R. Courant \& D. Hilbert - Methods of mathematical physics. Vol. II: Partial differential equations , Interscience Publishers, New York-London, 1962
\bibitem{CyFrKiSi} H. L. Cycon, R. G. Froese, W. Kirsch \& B. Simon - Schrödinger operators with application to quantum mechanics and global geometry , Texts and Monographs in Physics , Springer-Verlag, Berlin, 1987
\bibitem{DGN} D. Danielli, N. Garofalo \& D.-M. Nhieu - Non-doubling Ahlfors measures, perimeter measures, and the characterization of the trace spaces of Sobolev functions in Carnot-Carathéodory spaces , Mem. Amer. Math. Soc. , vol. 182, no. 857 , American Mathematical Society, Providence, RI, 2006
\bibitem{DGN2} D. Danielli, N. Garofalo \& D.-M. Nhieu - “A partial solution of the isoperimetric problem for the Heisenberg group” , Forum Math. 20 (2008) no. 1, p. 99-143
\bibitem{DFV} N. De Ponti, S. Farinelli \& I. Y. Violo - “Pleijel nodal domain theorem in non-smooth setting” , Trans. Amer. Math. Soc. Ser. B 11 (2024), p. 1138-1182
\bibitem{De} M. Derridj - “Un problème aux limites pour une classe d’opérateurs du second ordre hypoelliptiques” , Ann. Inst. Fourier (Grenoble) 21 (1971) no. 4, p. 99-148
\bibitem{EL} S. Eswarathasan \& C. Letrouit - “Nodal sets of eigenfunctions of sub-Laplacians” , Internat. Math. Res. Notices (2023) no. 23, p. 20670-20700
\bibitem{Fo0} G. B. Folland - “A fundamental solution for a subelliptic operator” , Bull. Amer. Math. Soc. 79 (1973), p. 373-376
\bibitem{Fo} G. B. Folland - “Subelliptic estimates and function spaces on nilpotent Lie groups” , Ark. Mat. 13 (1975) no. 2, p. 161-207
\bibitem{FS} G. B. Folland \& E. M. Stein - “Estimates for the ∂ ¯ b complex and analysis on the Heisenberg group” , Comm. Pure Appl. Math. 27 (1974), p. 429-522
\bibitem{FrSeSC} B. Franchi, R. Serapioni \& F. Serra Cassano - “Meyers-Serrin type theorems and relaxation of variational integrals depending on vector fields” , Houston J. Math. 22 (1996) no. 4, p. 859-890
\bibitem{FGL} R. L. Frank, D. Gontier \& M. Lewin - “The nonlinear Schrödinger equation for orthonormal functions II: Application to Lieb-Thirring inequalities” , Comm. Math. Phys. 384 (2021) no. 3, p. 1783-1828
\bibitem{FrHe1} R. L. Frank \& B. Helffer - “On Courant and Pleijel theorems for sub-Riemannian Laplacians” , 2024
\bibitem{FrHe2} R. L. Frank \& B. Helffer - “On Courant and Pleijel theorems for sub-Riemannian Laplacians” , in Pseudo-differential operators and related topics , Trends Math. , Birkhäuser/Springer, Cham, 2025, p. 9-23
\bibitem{FLW} R. L. Frank, A. Laptev \& T. Weidl - Schrödinger operators: eigenvalues and Lieb-Thirring inequalities , Cambridge Studies in Advanced Math. , vol. 200 , Cambridge University Press, Cambridge, 2023
\bibitem{FrLi} R. L. Frank \& E. H. Lieb - “Sharp constants in several inequalities on the Heisenberg group” , Ann. of Math. (2) 176 (2012) no. 1, p. 349-381
\bibitem{GaNh} N. Garofalo \& D.-M. Nhieu - “Isoperimetric and Sobolev inequalities for Carnot-Carathéodory spaces and the existence of minimal surfaces” , Comm. Pure Appl. Math. 49 (1996) no. 10, p. 1081-1144
\bibitem{Go} R. W. Goodman - Nilpotent Lie groups: structure and applications to analysis , Lect. Notes in Math. , vol. 562 , Springer-Verlag, Berlin-New York, 1976
\bibitem{HaKo} P. Hajłasz \& P. Koskela - Sobolev met Poincaré , Mem. Amer. Math. Soc. , vol. 145, no. 688 , American Mathematical Society, Providence, RI, 2000
\bibitem{HL} A. M. Hansson \& A. Laptev - “Sharp spectral inequalities for the Heisenberg Laplacian” , in Groups and analysis , London Math. Soc. Lecture Note Ser. , vol. 354 , Cambridge University Press, Cambridge, 2008, p. 100-115
\bibitem{HS} A. Hassannezhad \& D. Sher - “On Pleijel’s nodal domain theorem for the Robin problem” , Bull. London Math. Soc. 56 (2024) no. 4, p. 1449-1467
\bibitem{HeKiMa} J. Heinonen, T. Kilpeläinen \& O. Martio - Nonlinear potential theory of degenerate elliptic equations , Dover Publications, Inc., Mineola, NY, 2006, Unabridged republication of the 1993 original
\bibitem{HN} B. Helffer \& J. F. Nourrigat - “Approximation d’un système de champs de vecteurs et applications à l’hypoellipticité” , Ark. Mat. 17 (1979) no. 2, p. 237-254
\bibitem{Hel} B. Helffer - “Conditions nécessaires d’hypoanalyticité pour des opérateurs invariants à gauche homogènes sur un groupe nilpotent gradué” , J. Differential Equations 44 (1982) no. 3, p. 460-481
\bibitem{HN1} B. Helffer \& J. Nourrigat - Hypoellipticité maximale pour des opérateurs polynômes de champs de vecteurs , Progress in Math. , vol. 58 , Birkhäuser Boston, Inc., Boston, MA, 1985
\bibitem{HPS} B. Helffer \& M. Persson Sundqvist - “On nodal domains in Euclidean balls” , Proc. Amer. Math. Soc. 144 (2016) no. 11, p. 4777-4791
\bibitem{HeNoSw} J. G. Heywood, E. S. Noussair \& C. A. Swanson - “On the zeros of solutions of elliptic inequalities in bounded domains” , J. Differential Equations 28 (1978) no. 3, p. 345-353
\bibitem{Ho} L. Hörmander - “Hypoelliptic second order differential equations” , Acta Math. 119 (1967), p. 147-171
\bibitem{Je} F. Jean - Control of nonholonomic systems: from sub-Riemannian geometry to motion planning , SpringerBriefs in Math. , Springer, Cham, 2014
\bibitem{JL} D. Jerison \& J. M. Lee - “Extremals for the Sobolev inequality on the Heisenberg group and the CR Yamabe problem” , J. Amer. Math. Soc. 1 (1988) no. 1, p. 1-13
\bibitem{Jer} D. S. Jerison - “The Dirichlet problem for the Kohn Laplacian on the Heisenberg group. I” , J. Functional Analysis 43 (1981) no. 1, p. 97-142
\bibitem{LaWe} A. Laptev \& T. Weidl - “Sharp Lieb-Thirring inequalities in high dimensions” , Acta Math. 184 (2000) no. 1, p. 87-111
\bibitem{Le} C. Léna - “Pleijel’s nodal domain theorem for Neumann and Robin eigenfunctions” , Ann. Inst. Fourier (Grenoble) 69 (2019) no. 1, p. 283-301
\bibitem{LeRi} G. P. Leonardi \& S. Rigot - “Isoperimetric sets on Carnot groups” , Houston J. Math. 29 (2003) no. 3, p. 609-637
\bibitem{LeMa} G. P. Leonardi \& S. Masnou - “On the isoperimetric problem in the Heisenberg group $\mathbb{H}$ n ” , Ann. Mat. Pura Appl. (4) 184 (2005) no. 4, p. 533-553
\bibitem{Leo} G. Leoni - A first course in Sobolev spaces , Graduate Studies in Math. , vol. 181 , American Mathematical Society, Providence, RI, 2017
\bibitem{LMP} M. Levitin, D. Mangoubi \& I. Polterovich - Topics in spectral geometry , Graduate Studies in Mathematics , vol. 237 , American Mathematical Society, Providence, RI, 2023
\bibitem{LiLo} E. H. Lieb \& M. Loss - Analysis , Graduate Studies in Math. , vol. 14 , American Mathematical Society, Providence, RI, 2001
\bibitem{LiTh} E. H. Lieb \& W. E. Thirring - “Inequalities for the moments of the eigenvalues of the Schrödinger Hamiltonian and their relation to Sobolev inequalities” , Stud. math. Phys. (1976), p. 269-303, Essays honor Valentine Bargmann
\bibitem{Ma} F. Maggi - Sets of finite perimeter and geometric variational problems. An introduction to geometric measure theory , Cambridge Studies in Advanced Math. , vol. 135 , Cambridge University Press, Cambridge, 2012
\bibitem{MOS} W. Magnus, F. Oberhettinger \& R. P. Soni - Formulas and theorems for the special functions of mathematical physics , Grundlehren Math. Wissen. , vol. 52 , Springer-Verlag New York, Inc., New York, 1966
\bibitem{MS1} A. Menikoff \& J. Sjöstrand - “On the eigenvalues of a class of hypoelliptic operators” , Math. Ann. 235 (1978) no. 1, p. 55-85
\bibitem{MS2} A. Menikoff \& J. Sjöstrand - “On the eigenvalues of a class of hypoelliptic operators. II” , in Global analysis (Proc. Biennial Sem. Canad. Math. Congr., Univ. Calgary, Calgary, Alta., 1978) , Lect. Notes in Math. , vol. 755 , Springer, Berlin, 1979, p. 201-247
\bibitem{Me} G. Métivier - “Fonction spectrale et valeurs propres d’une classe d’opérateurs non elliptiques” , Comm. Partial Differential Equations 1 (1976) no. 5, p. 467-519
\bibitem{Me2} G. Métivier - “Hypoellipticité analytique sur des groupes nilpotents de rang 2 ” , Duke Math. J. 47 (1980) no. 1, p. 195-221
\bibitem{Miln} J. W. Milnor - Topology from the differentiable viewpoint , Princeton Landmarks in Math. , Princeton University Press, Princeton, NJ, 1997
\bibitem{Mo} A. Mohamed - “Étude spectrale d’opérateurs hypoelliptiques à caractéristiques multiples. I” , Ann. Inst. Fourier (Grenoble) 32 (1982) no. 3, p. 39-90
\bibitem{Mon} R. Monti - “Heisenberg isoperimetric problem. The axial case” , Adv. Calc. Var. 1 (2008) no. 1, p. 93-121
\bibitem{MM} R. Monti \& D. Morbidelli - “Non-tangentially accessible domains for vector fields” , Indiana Univ. Math. J. 54 (2005) no. 2, p. 473-498
\bibitem{MoRi} R. Monti \& M. Rickly - “Convex isoperimetric sets in the Heisenberg group” , Ann. Scuola Norm. Sup. Pisa Cl. Sci. (5) 8 (2009) no. 2, p. 391-415
\bibitem{MP} E. Müller-Pfeiffer - “On the number of nodal domains for eigenfunctions of elliptic differential operators” , J. London Math. Soc. (2) 31 (1985) no. 1, p. 91-100
\bibitem{Pa2} P. Pansu - “An isoperimetric inequality on the Heisenberg group” , Rend. Sem. Mat. Univ. e Politec. Torino (1983), p. 159-174, Conference on differential geometry on homogeneous spaces (Turin, 1983)
\bibitem{Pa} P. Pansu - “Une inégalité isopérimétrique sur le groupe de Heisenberg” , C. R. Acad. Sci. Paris Sér. I Math. 295 (1982) no. 2, p. 127-130
\bibitem{Pl} Å. Pleijel - “Remarks on Courant’s nodal line theorem” , Comm. Pure Appl. Math. 9 (1956), p. 543-550
\bibitem{Po} I. Polterovich - “Pleijel’s nodal domain theorem for free membranes” , Proc. Amer. Math. Soc. 137 (2009) no. 3, p. 1021-1024
\bibitem{Ri} M. Ritoré - “A proof by calibration of an isoperimetric inequality in the Heisenberg group $\mathbb{H}$ n ” , Calc. Var. Partial Differential Equations 44 (2012) no. 1-2, p. 47-60
\bibitem{RR} M. Ritoré \& C. Rosales - “Rotationally invariant hypersurfaces with constant mean curvature in the Heisenberg group $\mathbb{H}$ n ” , J. Geom. Anal. 16 (2006) no. 4, p. 703-720
\bibitem{RR2} M. Ritoré \& C. Rosales - “Area-stationary surfaces in the Heisenberg group $\mathbb{H}$ 1 ” , Adv. Math. 219 (2008) no. 2, p. 633-671
\bibitem{Rot} L. P. Rothschild - “A criterion for hypoellipticity of operators constructed from vector fields” , Comm. Partial Differential Equations 4 (1979) no. 6, p. 645-699
\bibitem{RS} L. P. Rothschild \& E. M. Stein - “Hypoelliptic differential operators and nilpotent groups” , Acta Math. 137 (1976) no. 3-4, p. 247-320
\bibitem{Str2} R. S. Strichartz - “Sub-Riemannian geometry” , J. Differential Geom. 24 (1986) no. 2, p. 221-263
\bibitem{St} R. S. Strichartz - “Estimates for sums of eigenvalues for domains in homogeneous spaces” , J. Functional Analysis 137 (1996) no. 1, p. 152-190
\bibitem{Na} B. v. Sz. Nagy - “Über Integralungleichungen zwischen einer Funktion und ihrer Ableitung” , Acta Sci. Math. (Szeged) 10 (1941), p. 64-74
\bibitem{Ta} G. Talenti - “The standard isoperimetric theorem” , in Handbook of convex geometry, Vol. A, B , North-Holland, Amsterdam, 1993, p. 73-123
\bibitem{Va} N. T. Varopoulos - “Analysis on nilpotent groups” , J. Functional Analysis 66 (1986) no. 3, p. 406-431
\bibitem{Wa} K. Watanabe - “Sur l’unicité du prolongement des solutions des équations elliptiques dégénérées” , Tohoku Math. J. (2) 34 (1982) no. 2, p. 239-249
